\begin{table}[!htbp]
\centering
\caption{Bargaining-Power Robustness: Model Contribution to Empirical Time Series in 2014}
\label{table-9-bargaining-power-robustness}
\resizebox{0.88\textwidth}{!}{%
\begin{tabular}{lrrr}
\toprule
Indicator & Baseline & $\zeta=0.4$ & $\zeta=0.6$ \\
\midrule
Unemployment & 32.9\% & 33.2\% & 32.8\% \\
Private Savings Rate & 20.7\% & 27.0\% & 16.7\% \\
Net Foreign Assets/GDP & 13.0\% & 12.1\% & 13.9\% \\
Current Account/GDP & 9.2\% & 8.9\% & 9.5\% \\
Current Account vis-a-vis Eurozone/GDP & 44.9\% & 43.7\% & 46.5\% \\
Unit Labor Costs Relative to Eurozone & 189.2\% & 284.2\% & 128.0\% \\
Private Savings Rate Rest of Eurozone & 4.1\% & 3.9\% & 4.3\% \\
\bottomrule
\end{tabular}}
\begin{flushleft}\footnotesize \textit{Notes:} Each entry reports the share of the observed empirical movement explained by the model-implied average change in calendar year 2014. The robustness columns change worker bargaining power to the displayed value while keeping the reform cut and all other baseline settings fixed.\end{flushleft}
\end{table}
